\pdfinfoomitdate=1
\pdftrailerid{}
\pdfsuppressptexinfo=15
\documentclass[11pt]{article}
\usepackage[margin=25mm]{geometry}
\usepackage{booktabs}
\usepackage[T1]{fontenc}
\usepackage[utf8]{inputenc}
\pagestyle{plain}
\title{Adaptive Prism Sensor Calibration: Version 1}
\author{LazyingArt Synthetic Evidence Lab}
\date{Project-owned synthetic sample, 2026-09-04}

\begin{document}
\maketitle
\begin{center}
\fbox{\parbox{0.9\linewidth}{\textbf{Evidence boundary:} project-owned synthetic source; not a benchmark, customer result, scientific result, or paid delivery.}}
\end{center}

\section*{Abstract}
This synthetic study describes a prism-lattice sensor and its amber drift
coefficient. Thermal stabilization precedes every measurement. The record is
deliberately small so that extraction, version identity, retrieval, and page
citations can be inspected by hand.

\section*{Purpose}
Version 1 establishes a local lexical baseline. It does not claim that the
fictional instrument exists or that the method generalizes to real papers.

\newpage
\section*{Calibration method}
The response model is
\[
  r(t) = g s(t) + b.
\]
The symbols may be linearized by a PDF text extractor, so the equation must be
checked against the source page rather than trusted as plain text alone.

\begin{center}
\begin{tabular}{ll}
\toprule
Setting & Version 1 value \\
\midrule
Calibration interval & 18 minutes \\
Sampling rate & 240 Hz \\
Reference wavelength & 532 nm \\
\bottomrule
\end{tabular}
\end{center}

The table is intentionally simple, but extracted reading order is still an
explicit weakness of the plain-text baseline.

\newpage
\section*{Limits and decision}
The declared temperature range is 18 to 24 degrees Celsius. Humidity coupling
was not tested. Local calibration reduces spectral drift only inside this
fictional setup.

Version 1 is suitable for a bounded extraction and citation demonstration. It
is not evidence for OCR, semantic retrieval, a knowledge graph, or production
deployment.
\end{document}
